\begin{textsample}{NEG Dim 4 – vem\_ed\_24 – Score 1.00 – vem\_ed\_24/t130.txt}  \label{ex:f4_neg_063}
Artigo de opinião Do Teletandem aos PCIs / Cesu: colaboração para a internacionalização em casa Priscilla Ferro priscilla.ferro@cps.sp.gov.br No contexto atual do ensino superior, os Intercâmbios Virtuais (Virtual Exchange) têm ganhado papel de protagonismo (O'Dowd, 2018).
Trata-se de práticas educacionais que utilizam tecnologias de comunicação para facilitar as interações entre indivíduos ou grupos geograficamente separados e, dessa forma, promovem a aprendizagem intercultural e o desenvolvimento de competências linguísticas e tecnológicas.
Há várias abordagens e metodologias para os Intercâmbios Virtuais, duas das abordagens mais conhecidas são o Teletandem (Telles e Vassallo, 2006) e o Collaborative Online International Learning (COIL), como são mundialmente conhecidos nossos PCIs (Projetos Colaborativos Internacionais).
Ambas têm os mesmos objetivos de expandir fronteiras educacionais e promover os intercâmbios cultural e linguístico entre estudantes de diferentes partes do mundo.
O Teletandem, formalizado em 2006 na Universidade Estadual Paulista (Unesp), evoluiu das práticas de tandem de aprendizagem de línguas estrangeiras (Telles e Vassallo, 2006).
Estudantes de diferentes países utilizam mensageiros instantâneos como o Skype e se conectam para aprender a língua um do outro, promovendo um aprendizado recíproco e enriquecedor.
Com a visão de"eu ajudo você a aprender a minha língua e você me ajuda a aprender a sua", o Teletandem é um ambiente de aprendizagem de línguas estrangeiras que vem se estabelecendo com sólidas bases de pesquisas e de práticas.
Os projetos COIL, por sua vez, ampliam a colaboração para além das línguas estrangeiras, pois abraçam disciplinas variadas (Guth e Helm, 2010).
Essa prática integra a colaboração internacional nos currículos acadêmicos e permite que estudantes, de diversas áreas do saber, trabalhem colaborativamente.
Nesses trabalhos, é possível haver diferentes focos, como, por exemplo, em competências linguísticas e interculturais ou em conteúdos interdisciplinares e transdisciplinares e, seja qual for, sua adoção no ensino superior traz benefícios significativos para alunos e professores. continuação Para os alunos, a aprendizagem prática e contextualizada oferece a oportunidade de aplicar conhecimentos teóricos em contextos reais, em parceria com colegas de diferentes países em projetos que refletem situações do mundo real (De Wit, 2011).
A interação com diferentes culturas desenvolve competências globais, como compreensão e empatia intercultural, habilidades essenciais no mercado de trabalho globalizado.
Soft skills como comunicação, liderança, resolução de problemas e trabalho em equipe também são aprimoradas.
Para os professores, os benefícios incluem novas perspectivas e métodos de ensino, o que enriquece o conteúdo acadêmico e torna as aulas mais dinâmicas.
A implementação de COIL promove a inovação no ensino, incentiva o uso de tecnologias educacionais e de metodologias ativas de aprendizagem (Rubin e Guth, 2023).
Professores envolvidos em projetos COIL têm a oportunidade de ampliar suas redes de contatos internacionais e de participar de pesquisas colaborativas, visando a publicações e, eventualmente, podem ministrar aulas em clases espejo, quando o professor de uma das instituições ministra uma aula para a turma da instituição parceira (o que rende um certificado de aula internacional).
No entanto, há professores de disciplinas técnicas ou profissionais que se sentem limitados pela barreira da língua estrangeira.
Para esses casos, faz-se necessário o apoio dos professores de línguas estrangeiras aos docentes de conteúdos específicos (Belz, 2003).
Os professores de línguas auxiliam na superação dessas barreiras linguísticas, garantindo uma comunicação eficaz para o desenvolvimento do projeto, além de trabalharem nuances culturais e práticas de comunicação intercultural, essenciais para o sucesso de colaborações internacionais.
Professores da Fatec que embarcam nos PCIs têm, portanto, uma grande vantagem, já que a instituição conta com mais de 280 professores de língua inglesa e mais de 80 de língua espanhola.
Esse corpo docente pode oferecer apoio focado em vocabulário específico e acompanhamento sobre a comunicação intercultural.
A colaboração entre professores de línguas e de conteúdos específicos resulta em projetos mais ricos e integrados, nos quais as aprendizagens de línguas e de disciplinas ocorrem de forma harmoniosa (Thorne, 2006).
Essa sinergia é um dos ingredientes para o sucesso das iniciativas de internacionalização, pois enriquece a experiência educativa e prepara melhor os alunos para atuarem em um mundo globalizado.
A integração de professores de línguas e de outras disciplinas nos projetos COIL representa um avanço significativo na internacionalização da educação.
Ao alinhar os esforços desses professores, cria-se uma educação mais inclusiva, eficaz e globalmente conectada.
Os Intercâmbios Virtuais são excelentes maneiras de compor a preparação dos estudantes para atuar em um mundo interdependente e em constante transformação e de se construir um futuro no qual o aprendizado transcenda fronteiras e promova a paz e a compreensão global por meio da educação.
Referências BELZ, Julie A.
Linguistic Perspectives on the Development of Intercultural Competence in Telecollaboration.
Language, Learning \& Technology.
Honolulu, v.7, n.2, p. 68-99, May 2003.
Disponível em: https://scholarspace.manoa.hawaii.edu/server/api/core/bitstreams/dec7d29d-63a7-4715-823f-66ae2e8021dd/content.
Acesso em 26 \textbf{jul}. 2024.
DE WIT, Hans.
Internationalisation of Higher Education: Nine Misconceptions.
International Higher Education, n.64, p. 1-6, Summer 2011.
Disponível em: https://ejournals.bc.edu/index.php/ihe/article/view/8556/8321.
Acesso em: 26 \textbf{jul}. 2024 GUTH, Sarah; HELM, Francesca.
Telecollaboration 2.0: Language, Literacies and Intercultural Learning in the 21st Century.
Bern: Peter Lang, 2010.
O'DOWD, Robert.
From telecollaboration to virtual exchange: state-of-the-art and the role of unicollaboration in moving forward.
Journal Of Virtual Exchange, [ S.L. ], v. 1, p. 1-23, 24 abr. 2018.
Disponível em: https://journal.unicollaboration.org/article/view/35567/33147.
Acesso em: 25 \textbf{jul}. 2018.
RUBIN, Jon; GUTH, Sarah.
The Guide to COIL Virtual Exchange.
New York: Routledge, 2023.
TELLES, João; VASSALO, Maria Luisa.
Foreign language learning In-Tandem: Teletandem as an alternative proposal in CALLT.
The ESPecialist, São Paulo, v. 27, n.2., p. 189-212, 2006.
Disponível em: https://revistas.pucsp.br/index.php/esp/article/view/1629/1048.
Acesso em: 25 \textbf{jul}. 2024.
THORNE, Steven L.
Pedagogical and praxiological lessons from internet-mediated intercultural foreign language education research.
In: Internet-mediated intercultural foreign language education.
Boston: Heinle, Cengage Learning, 2006.
Disponível em: https://scholarspace.manoa.hawaii.edu/server/api/core/bitstreams/b0445d68-2405-444b-b667-4ae3bf198c4a/content.
Acesso em: 25 jul.2024.

% matched lemmas: jul
\end{textsample}
